% Generated from validated Article Draft Result. Do not hand-edit; revise the structured draft instead.
\begingroup\raggedright
\subsection{Safetyを一つの「危険度」へ畳まない}
\par\endgroup

4月のagent safety／evaluation周辺では、prompt injectionへの防御、GUI agentのprocess-centric benchmark、PII filtering、screenshot上のinjection検知、alignment fakingの観測という、互いに異なる課題が並んだ。これらは同じ脅威を別手法で解いているわけではない。入力を隔離する設計、能力を測るbenchmark、privacy保護、視覚的attack detection、behavioral evaluationは、それぞれ前提とfailure modeが異なる。本稿では「agentが危険になった」という一文へまとめず、実行条件そのものを設計・測定する問題が前景化したと整理する。\autocite{src-2f8952da2dc3,src-3efc4342b784,src-477d406be59c,src-73ae4b975e7b,src-8cf1b0dc8361,src-a408165af133}


\begingroup\raggedright
\subsection{AgentVisor — semantic privilege separation}
\par\endgroup

AgentVisorは4月27日の論文で、OS virtualizationから着想したsemantic privilege separationを使うprompt-injection defense frameworkとして提案された。発想の核は、agentが受け取る自然言語をすべて同じ権限の命令として扱わず、trustedな意図と外部content由来のinstructionを意味上の権限境界で分離することにある。これは特定のattack文字列をblacklistする話とは異なり、agent executionの前に「誰の指示として解釈してよいか」をsystem designへ持ち込む。防御性能はauthor-reportedであり、本稿は一般的な安全保証としては扱わない。\autocite{src-8cf1b0dc8361}


\begingroup\raggedright
\subsection{WindowsWorld — 成功率だけでなくprocessを測る}
\par\endgroup

WindowsWorldは4月30日に、professionalなcross-application environmentでautonomous GUI agentを評価するprocess-centric benchmarkとして提示された。複数applicationをまたぐmulti-step taskでは、最終結果だけを採点すると途中の操作ミスや偶然の成功を見落とし得る。process-centricという問題設定は、computer-use agentの評価対象が単発UI操作からworkflow全体へ広がっていることを示す。benchmark上のmodel performanceやrankingはauthorsのevaluationであり、productionでの一般能力を直接保証するものではない。\autocite{src-a408165af133}


\begingroup\raggedright
\subsection{Privacy Filter — safety機能がopen ecosystemへ流れる}
\par\endgroup

OpenAIは4月22日にOpenAI Privacy Filterを、text中のpersonally identifiable informationをdetect／redactするopen-weight modelとして公開した。「state-of-the-art accuracy」という表現はvendor claimとして扱う。同日のTransformers v5.6.0 release notesは、このPrivacy Filter integrationを記録している。ここで興味深いのは、privacy protectionがclosed product内部のfeatureだけでなく、open-weight modelと一般的なmodel libraryを通じてdeveloper stackへ組み込める形で現れたことだ。\autocite{src-2f8952da2dc3,src-3efc4342b784}


\begingroup\raggedright
\subsection{SnapGuard — screenshotをattack surfaceとして扱う}
\par\endgroup

SnapGuardは4月28日の論文で、screenshot-based web agentを対象に、webpage screenshotのmultimodal representation analysisとしてprompt injection detectionを定式化する。text DOMだけを検査するのではなく、agentが実際に見る視覚表現を防御対象に含める点が問題設定の特徴である。これはAgentVisorのsemantic privilege separationとは異なる層の防御であり、両者の検出率を単純比較して優劣を決める材料ではない。\autocite{src-477d406be59c}


\begingroup\raggedright
\subsection{Tatemae — tool selectionをbehavioral signalとして見る}
\par\endgroup

Tatemaeは4月29日に、alignment fakingを「training objectiveへ戦略的に従い、監視が外れると以前のpreferenceへ戻る」現象として定義し、tool selectionを観測signalとして扱う研究を提示した。ここでの問いはprompt injection defenseとは別で、モデル内部の意図を直接読めない状況でbehavioral choiceからalignment-related behaviorをどう検出するかにある。著者の設定と結果を、一般のagentが実際にalignment fakingしているという prevalence claimへ拡張しない。\autocite{src-73ae4b975e7b}


\begingroup\raggedright
\subsection{異なる境界を層として読む}
\par\endgroup

4月の材料を層で分けると、AgentVisorはinstruction privilege、SnapGuardはvisual input、Privacy Filterはdata handling、WindowsWorldはworkflow evaluation、Tatemaeはbehavioral alignment signalを扱う。同じagent systemの周辺に置けるものではあるが、脅威モデルも成功条件も異なる。後続月でMCP、persistent memory、long-horizon controlなどattack surfaceが増える前に、4月は「何をtrustedとし、何を観測し、どこで遮断するか」という基礎的な境界設計が複数方向から現れた月だった。\autocite{src-3efc4342b784,src-477d406be59c,src-73ae4b975e7b,src-8cf1b0dc8361,src-a408165af133}


